\cleardoublepage
\selectlanguage{spanish}
\chapter*{Resumen}
\addcontentsline{toc}{chapter}{Resumen}

El análisis y la cualificación de licitaciones públicas requiere revisar
documentación técnica y administrativa de gran extensión para identificar los
requisitos exigidos, evaluar la solvencia de la empresa licitadora y valorar la
viabilidad del contrato. Cuando este trabajo se realiza manualmente, supone una
carga significativa para los equipos responsables y puede dar lugar a errores
de omisión.

Este Trabajo Fin de Grado aborda el diseño y desarrollo de LIA, una plataforma
web de gestión de oportunidades comerciales centrada en licitaciones públicas.
La solución permite organizar la cartera de oportunidades mediante pipelines
configurables y flujos de trabajo estructurados, e integra capacidades de
inteligencia artificial generativa para apoyar la extracción de información
clave, la evaluación preliminar de solvencia y la generación de documentos de
análisis.

El objetivo general del proyecto es aligerar el proceso de análisis y toma de
requisitos de las licitaciones públicas. Para ello, se plantea automatizar la
detección de solvencia económica, optimizar la relación entre coste y
rendimiento de los proveedores de inteligencia artificial generativa y
facilitar la creación de flujos de trabajo adaptados al análisis de
licitaciones. La plataforma se desarrolla como un monorepositorio con
Turborepo, una aplicación web basada en Vue 3 y una API implementada con
NestJS.

\vspace{0.5cm}
\noindent\textbf{Palabras clave:} licitaciones públicas, inteligencia
artificial generativa, gestión de oportunidades, análisis de solvencia,
flujos de trabajo.
